\documentclass[10pt,a4paper]{amsart}
\usepackage{amsmath,amssymb,amsthm,graphicx}
\usepackage{verbatim}
\newenvironment{reference}{\comment}{\endcomment}
%\newenvironment{reference}{}{}
\newenvironment{slogan}{\comment}{\endcomment}
\newenvironment{history}{\comment}{\endcomment}

% For commutative diagrams we use Xy-pic
\usepackage[all]{xy}

% We use 2cell for 2-commutative diagrams.
\xyoption{2cell}
% Theorem environments.
%
\theoremstyle{plain}
\newtheorem{theorem}[subsection]{Theorem}
\newtheorem{proposition}[subsection]{Proposition}
\newtheorem{lemma}[subsection]{Lemma}

\theoremstyle{definition}
\newtheorem{definition}[subsection]{Definition}
\newtheorem{example}[subsection]{Example}
\newtheorem{exercise}[subsection]{Exercise}
\newtheorem{situation}[subsection]{Situation}

\theoremstyle{remark}
\newtheorem{remark}[subsection]{Remark}
\newtheorem{remarks}[subsection]{Remarks}

\numberwithin{equation}{subsection}

% Macros
%
\def\lim{\mathop{\mathrm{lim}}\nolimits}
\def\colim{\mathop{\mathrm{colim}}\nolimits}
\def\Spec{\mathop{\mathrm{Spec}}}
\def\Hom{\mathop{\mathrm{Hom}}\nolimits}
\def\Ext{\mathop{\mathrm{Ext}}\nolimits}
\def\SheafHom{\mathop{\mathcal{H}\!\mathit{om}}\nolimits}
\def\SheafExt{\mathop{\mathcal{E}\!\mathit{xt}}\nolimits}
\def\Sch{\mathit{Sch}}
\def\Mor{\mathop{\mathrm{Mor}}\nolimits}
\def\Ob{\mathop{\mathrm{Ob}}\nolimits}
\def\Sh{\mathop{\mathit{Sh}}\nolimits}
\def\NL{\mathop{N\!L}\nolimits}
\def\CH{\mathop{\mathrm{CH}}\nolimits}
\def\proetale{{pro\text{-}\acute{e}tale}}
\def\etale{{\acute{e}tale}}
\def\QCoh{\mathit{QCoh}}
\def\Ker{\mathop{\mathrm{Ker}}}
\def\Im{\mathop{\mathrm{Im}}}
\def\Coker{\mathop{\mathrm{Coker}}}
\def\Coim{\mathop{\mathrm{Coim}}}

% Boxtimes
%
\DeclareMathSymbol{\boxtimes}{\mathbin}{AMSa}{"02}

%
% Macros for moduli stacks/spaces
%
\def\QCohstack{\mathcal{QC}\!\mathit{oh}}
\def\Cohstack{\mathcal{C}\!\mathit{oh}}
\def\Spacesstack{\mathcal{S}\!\mathit{paces}}
\def\Quotfunctor{\mathrm{Quot}}
\def\Hilbfunctor{\mathrm{Hilb}}
\def\Curvesstack{\mathcal{C}\!\mathit{urves}}
\def\Polarizedstack{\mathcal{P}\!\mathit{olarized}}
\def\Complexesstack{\mathcal{C}\!\mathit{omplexes}}
% \Pic is the operator that assigns to X its picard group, usage \Pic(X)
% \Picardstack_{X/B} denotes the Picard stack of X over B
% \Picardfunctor_{X/B} denotes the Picard functor of X over B
\def\Pic{\mathop{\mathrm{Pic}}\nolimits}
\def\Picardstack{\mathcal{P}\!\mathit{ic}}
\def\Picardfunctor{\mathrm{Pic}}
\def\Deformationcategory{\mathcal{D}\!\mathit{ef}}
\usepackage{cleveref}
\crefname{originalsection}{original section}{original sections}
\Crefname{originalsection}{Original section}{Original sections}
\makeatletter
\newlabel{homology-example-graded-vector-spaces}{{0FFX}{source}{}{sourcepointer-1}{}}
\newlabel{homology-example-graded-vector-spaces@cref}{{[originalsection][0][]0FFX}{[1][0][]source}{}{}{}}
\makeatother
\crefname{originallemma}{original lemma}{original lemmas}
\Crefname{originallemma}{Original lemma}{Original lemmas}
\makeatletter
\newlabel{homology-lemma-left-dual-graded-vector-spaces}{{0FFV}{source}{}{sourcepointer-2}{}}
\newlabel{homology-lemma-left-dual-graded-vector-spaces@cref}{{[originallemma][0][]0FFV}{[1][0][]source}{}{}{}}
\makeatother
\crefname{originalreference}{original reference}{original references}
\Crefname{originalreference}{Original reference}{Original references}
\makeatletter
\newlabel{lemma-composition-correspondences}{{0FG0}{source}{}{sourcepointer-3}{}}
\newlabel{lemma-composition-correspondences@cref}{{[originalreference][0][]0FG0}{[1][0][]source}{}{}{}}
\makeatother
\crefname{originalreference}{original reference}{original references}
\Crefname{originalreference}{Original reference}{Original references}
\makeatletter
\newlabel{lemma-dual}{{08I1}{source}{}{sourcepointer-4}{}}
\newlabel{lemma-dual@cref}{{[originalreference][0][]08I1}{[1][0][]source}{}{}{}}
\makeatother
\crefname{originalreference}{original reference}{original references}
\Crefname{originalreference}{Original reference}{Original references}
\makeatletter
\newlabel{lemma-from-functor-to-weil-classical}{{0FH1}{source}{}{sourcepointer-5}{}}
\newlabel{lemma-from-functor-to-weil-classical@cref}{{[originalreference][0][]0FH1}{[1][0][]source}{}{}{}}
\makeatother
\crefname{originalreference}{original reference}{original references}
\Crefname{originalreference}{Original reference}{Original references}
\makeatletter
\newlabel{lemma-functor-and-cycles}{{0FG5}{source}{}{sourcepointer-6}{}}
\newlabel{lemma-functor-and-cycles@cref}{{[originalreference][0][]0FG5}{[1][0][]source}{}{}{}}
\makeatother
\crefname{originalreference}{original reference}{original references}
\Crefname{originalreference}{Original reference}{Original references}
\makeatletter
\newlabel{lemma-prep-dual}{{0FG8}{source}{}{sourcepointer-7}{}}
\newlabel{lemma-prep-dual@cref}{{[originalreference][0][]0FG8}{[1][0][]source}{}{}{}}
\makeatother
\crefname{originalreference}{original reference}{original references}
\Crefname{originalreference}{Original reference}{Original references}
\makeatletter
\newlabel{remark-replace-cup-product-classical}{{0FGU}{source}{}{sourcepointer-8}{}}
\newlabel{remark-replace-cup-product-classical@cref}{{[originalreference][0][]0FGU}{[1][0][]source}{}{}{}}
\makeatother
\title{Producing the axioms of a Weil cohomology theory from a monoidal functor on motives}
\author{The Stacks Project Authors}
\date{stacks-project:weil.tex:local-snapshot-20261012; GNU FDL 1.2 or later}
\begin{document}
\maketitle
\noindent\textit{Editorial note.} All delivered mathematical text is the author's own, copied from the pinned Stacks Project chapter file weil.tex; no statement or proof step was rewritten, reordered or removed.\par\medskip
\noindent\textit{Editorial note.} The counted proof is the author's complete argument: The proof opens verbatim with 'This proof is the same as the proof of Lemma \textbackslash{}ref\{lemma-from-functor-to-weil-classical\}; we urge the reader to read the proof of that lemma instead.' — a reading suggestion, not an omission: the derivation is nevertheless written out in full in the span. Beyond that: the construction of the cohomology algebra, the identification of the coefficient object, the cycle class and degree maps from morphism groups of motives, and the axiom by axiom verification.\par\medskip
\noindent\textit{Editorial note.} The label lemma-from-functor-to-weil names the counted statement, and the proof environment immediately following it in the source is delivered verbatim as its proof.\par\medskip
\noindent\textit{Editorial note.} The counted statement is carried by the author's own lemma environment in the chapter source.\par\medskip
\noindent\textit{Editorial note.} Pre-open gate: python operations/agent-authoring/tools/precard\_gate.py --source weil.tex --label lemma-from-functor-to-weil --start 3009 --end 3194 --local-id 28134 returned RESULT=OK (L28134 free, occupancy FREE, public interval clear). Pre-assembly tag self-check: canonical tag resolved and not present in reference/public-items.jsonl.\par\medskip
\noindent\textit{Editorial note.} Editorial note (verbatim, LANE\_BRIEF §5 reading-suggestion rule): the proof opens with "This proof is the same as the proof of Lemma \textbackslash{}ref\{lemma-from-functor-to-weil-classical\}; we urge the reader to read the proof of that lemma instead." The proof environment nevertheless contains the full derivation, so the card is counted.\par\medskip
\noindent\textit{Editorial note.} The statement and proof are copied from the pinned chapter file; the cited inputs are kept as external pointers with their canonical Stacks tags.\par\medskip
\noindent\textit{Allowed standard prerequisites (editorial).}\begin{itemize}
\item Correspondences as morphism groups between motives. Same chapter, the lemma identifying the groups of correspondences of degree r from the base to a scheme with the morphism groups between the corresponding motives and with the rational Chow groups. Identifies the inputs of the cycle class map and of the degree map with morphism groups in the category of motives, which is what allows applying the functor.
\item Duality of the twisted motive of a smooth projective scheme. Same chapter, the lemma that the twisted motive of an equidimensional smooth projective scheme is a left dual of its motive, and the computation of the corresponding evaluation as the class of the diagonal. Produces the nondegenerate pairings on cohomology and identifies them with cup product followed by the degree map.
\item The projection formula and the description of the cycle class maps through correspondences. Same chapter, the lemmas relating pushforwards of cycles to compositions of correspondences, together with the preparation lemma for the duality diagram. Verifies the compatibility of the cycle class map with pushforwards and the projection formula.
\end{itemize}\medskip
\ifdefined\hypertarget\hypertarget{sourcepointer-1}{}\fi
\noindent\textit{Source reference.} Original section 0FFX: sources/stacks-chow/tags.txt, source lines 19302--19302. Prerequisite pointer invoked at proof line(s) 3046; it names a standard fact of a sibling chapter that the counted argument uses, and every step of the counted proof built on it is retained verbatim in the delivered body.\par\medskip
\ifdefined\hypertarget\hypertarget{sourcepointer-2}{}\fi
\noindent\textit{Source reference.} Original lemma 0FFV: sources/stacks-chow/tags.txt, source lines 19300--19300. Prerequisite pointer invoked at proof line(s) 3105; it names a standard fact of a sibling chapter that the counted argument uses, and every step of the counted proof built on it is retained verbatim in the delivered body.\par\medskip
\ifdefined\hypertarget\hypertarget{sourcepointer-3}{}\fi
\noindent\textit{Source reference.} Original reference 0FG0: sources/stacks-weil/pinned-original-source.tex, source lines 441--441. Prerequisite pointer invoked at proof line(s) 3067; it names a standard fact of the same chapter that the counted argument uses, and every step of the counted proof built on it is retained verbatim in the delivered body.\par\medskip
\ifdefined\hypertarget\hypertarget{sourcepointer-4}{}\fi
\noindent\textit{Source reference.} Original reference 08I1: sources/stacks-weil/pinned-original-source.tex, source lines 1186--1186. Prerequisite pointer invoked at proof line(s) 3100; it names a standard fact of the same chapter that the counted argument uses, and every step of the counted proof built on it is retained verbatim in the delivered body.\par\medskip
\ifdefined\hypertarget\hypertarget{sourcepointer-5}{}\fi
\noindent\textit{Source reference.} Original reference 0FH1: sources/stacks-weil/pinned-original-source.tex, source lines 1858--1858. Prerequisite pointer invoked at proof line(s) 3023; it names a standard fact of the same chapter that the counted argument uses, and every step of the counted proof built on it is retained verbatim in the delivered body.\par\medskip
\ifdefined\hypertarget\hypertarget{sourcepointer-6}{}\fi
\noindent\textit{Source reference.} Original reference 0FG5: sources/stacks-weil/pinned-original-source.tex, source lines 661--661. Prerequisite pointer invoked at proof line(s) 3132; it names a standard fact of the same chapter that the counted argument uses, and every step of the counted proof built on it is retained verbatim in the delivered body.\par\medskip
\ifdefined\hypertarget\hypertarget{sourcepointer-7}{}\fi
\noindent\textit{Source reference.} Original reference 0FG8: sources/stacks-weil/pinned-original-source.tex, source lines 777--777. Prerequisite pointer invoked at proof line(s) 3147; it names a standard fact of the same chapter that the counted argument uses, and every step of the counted proof built on it is retained verbatim in the delivered body.\par\medskip
\ifdefined\hypertarget\hypertarget{sourcepointer-8}{}\fi
\noindent\textit{Source reference.} Original reference 0FGU: sources/stacks-weil/pinned-original-source.tex, source lines 1653--1653. Prerequisite pointer invoked at proof line(s) 3180; it names a standard fact of the same chapter that the counted argument uses, and every step of the counted proof built on it is retained verbatim in the delivered body.\par\medskip
\begin{lemma}
\label{lemma-from-functor-to-weil}
Let $k$ be a field. Let $F$ be a field of characteristic $0$.
Assume given a $\mathbf{Q}$-linear functor
$$
G : M_k \longrightarrow \text{graded }F\text{-vector spaces}
$$
of symmetric monoidal categories such that $G(\mathbf{1}(1))$
is nonzero only in degree $-2$. Then we obtain data (D0), (D1), (D2), and (D3)
satisfying all of (A), (B), and (C) above.
\end{lemma}

\begin{proof}
This proof is the same as the proof of
Lemma \ref{lemma-from-functor-to-weil-classical};
we urge the reader to read the proof of that lemma instead.

\medskip\noindent
We obtain a contravariant functor from the category of smooth
projective schemes over $k$ to the category of graded $F$-vector spaces
by setting $H^*(X) = G(h(X))$. By assumption we have a canonical
isomorphism
$$
H^*(X \times Y) = G(h(X \times Y)) = G(h(X) \otimes h(Y)) =
G(h(X)) \otimes G(h(Y)) = H^*(X) \otimes H^*(Y)
$$
compatible with pullbacks. Using pullback along the diagonal
$\Delta : X \to X \times X$ we obtain a canonical map
$$
H^*(X) \otimes H^*(X) = H^*(X \times X) \to H^*(X)
$$
of graded vector spaces compatible with pullbacks.
This defines a functorial graded $F$-algebra structure on
$H^*(X)$. Since $\Delta$ commutes with the commutativity
constraint $h(X) \otimes h(X) \to h(X) \otimes h(X)$ (switching the factors)
and since $G$ is a functor of symmetric monoidal categories (so compatible with
commutativity constraints), and by our convention in
Homology, Example \ref{homology-example-graded-vector-spaces}
we conclude that $H^*(X)$ is a graded
commutative algebra! Hence we get our datum (D1).

\medskip\noindent
Since $\mathbf{1}(1)$ is invertible in the category of motives
we see that $G(\mathbf{1}(1))$ is invertible in the category of
graded $F$-vector spaces. Thus $\sum_i \dim_F G^i(\mathbf{1}(1)) = 1$.
By assumption we only get something nonzero in degree $-2$.
Our datum (D0) is the vector space $F(1) = G^{-2}(\mathbf{1}(1))$.
Since $G$ is a symmetric monoidal functor we see that
$F(n) = G^{-2n}(\mathbf{1}(n))$ for all $n \in \mathbf{Z}$.
It follows that
$$
H^{2r}(X)(r) = G^{2r}(h(X)) \otimes G^{-2r}(\mathbf{1}(r)) =
G^0(h(X)(r))
$$
a formula we will frequently use below.

\medskip\noindent
Let $X$ be a smooth projective scheme over $k$. By
Lemma \ref{lemma-composition-correspondences} we have
$$
\CH^r(X) \otimes \mathbf{Q} = \text{Corr}^r(\Spec(k), X) =
\Hom(\mathbf{1}(-r), h(X)) = \Hom(\mathbf{1}, h(X)(r))
$$
Applying the functor $G$ this maps into
$\Hom(G(\mathbf{1}), G(h(X)(r)))$.
By taking the image of $1$ in $G^0(\mathbf{1}) = F$
into $G^0(h(X)(r)) = H^{2r}(X)(r)$ we obtain
$$
\gamma :
\CH^r(X) \otimes \mathbf{Q} \longrightarrow H^{2r}(X)(r)
$$
This is the datum (D2).

\medskip\noindent
Let $X$ be a nonempty smooth projective scheme over $k$
which is equidimensional of dimension $d$. By
Lemma \ref{lemma-composition-correspondences} we have
$$
\Mor(h(X)(d), \mathbf{1}) = \Mor((X, 1, d), (\Spec(k), 1, 0)) =
\text{Corr}^{-d}(X, \Spec(k)) = \CH_d(X)
$$
Thus the class of the cycle $[X]$ in $\CH_d(X)$ defines a morphism
$h(X)(d) \to \mathbf{1}$. Applying $G$ and taking degree $0$
parts we obtain
$$
H^{2d}(X)(d) = G^0(h(X)(d)) \longrightarrow G^0(\mathbf{1}) = F
$$
This map $\int_X : H^{2d}(X)(d) \to F$ is the datum (D3).

\medskip\noindent
Let $X$ be a smooth projective scheme over $k$ which is
nonempty and equidimensional of dimension $d$. By Lemma \ref{lemma-dual}
we know that $h(X)(d)$ is a left dual to $h(X)$. Hence
$G(h(X)(d)) = H^*(X) \otimes_F F(d)[2d]$
is a left dual to $H^*(X)$ in the category of graded $F$-vector spaces.
Here $[n]$ is the shift functor on graded vector spaces.
By Homology, Lemma \ref{homology-lemma-left-dual-graded-vector-spaces}
we find that $\sum_i \dim_F H^i(X) < \infty$ and that
$\epsilon : h(X)(d) \otimes h(X) \to \mathbf{1}$ produces
nondegenerate pairings $H^{2d - i}(X)(d) \otimes_F H^i(X) \to F$.
In the proof of Lemma \ref{lemma-dual} we have seen that
$\epsilon$ is given by $[\Delta]$ via the identifications
$$
\Hom(h(X)(d) \otimes h(X), \mathbf{1}) =
\text{Corr}^{-d}(X \times X, \Spec(k)) =
\CH_d(X \times X)
$$
Thus $\epsilon$ is the composition of $[X] : h(X)(d) \to \mathbf{1}$
and $h(\Delta)(d) : h(X)(d) \otimes h(X) \to h(X)(d)$. It follows
that the pairings above are given by cup product followed by
$\int_X$. This proves axiom (A).

\medskip\noindent
Axiom (B) follows from the assumption that $G$ is compatible
with tensor structures and our construction of the cup product above.

\medskip\noindent
Axiom (C). Our construction of $\gamma$ takes a cycle $\alpha$ on $X$,
interprets it a correspondence $a$ from $\Spec(k)$ to $X$ of some degree,
and then applies $G$. If $f : Y \to X$ is a morphism of nonempty
equidimensional smooth projective schemes over $k$, then
$f^!\alpha$ is the pushforward (!) of $\alpha$
by the correspondence $[\Gamma_f]$ from $X$ to $Y$, see
Lemma \ref{lemma-functor-and-cycles}. Hence
$f^!\alpha$ viewed as a correspondence from $\Spec(k)$ to $Y$
is equal to $a \circ [\Gamma_f]$, see
Lemma \ref{lemma-composition-correspondences}.
Since $G$ is a functor, we conclude
$\gamma$ is compatible with pullbacks, i.e., axiom (C)(a) holds.

\medskip\noindent
Let $f : Y \to X$ be a morphism of nonempty equidimensional
smooth projective schemes over $k$ and
let $\beta \in \CH^r(Y)$ be a cycle on $Y$. We have to show that
$$
\int_Y \gamma(\beta) \cup f^*c = \int_X \gamma(f_*\beta) \cup c
$$
for all $c \in H^*(X)$. Let $a, a^t, \eta_X, \eta_Y, [X], [Y]$
be as in Lemma \ref{lemma-prep-dual}.
Let $b$ be $\beta$ viewed as a correspondence from $\Spec(k)$ to $Y$
of degree $r$. Then $f_*\beta$ viewed as a correspondence from
$\Spec(k)$ to $X$ is equal to $a^t \circ b$, see
Lemmas \ref{lemma-functor-and-cycles} and
\ref{lemma-composition-correspondences}.
The displayed equality above holds if we can show that
$$
h(X) = \mathbf{1} \otimes h(X)
\xrightarrow{b \otimes 1}
h(Y)(r) \otimes h(X)
\xrightarrow{1 \otimes a}
h(Y)(r) \otimes h(Y)
\xrightarrow{\eta_Y}
h(Y)(r)
\xrightarrow{[Y]}
\mathbf{1}(r - e)
$$
is equal to
$$
h(X) = \mathbf{1} \otimes h(X)
\xrightarrow{a^t \circ b \otimes 1}
h(X)(r + d - e) \otimes h(X)
\xrightarrow{\eta_X}
h(X)(r + d - e)
\xrightarrow{[X]}
\mathbf{1}(r - e)
$$
This follows immediately from Lemma \ref{lemma-prep-dual}.
Thus we have axiom (C)(b).

\medskip\noindent
To prove axiom (C)(c) we use the discussion in
Remark \ref{remark-replace-cup-product-classical}.
Hence it suffices to prove that $\gamma$ is compatible with
exterior products. Let $X$, $Y$ be nonempty smooth projective
schemes over $k$ and let $\alpha$, $\beta$ be cycles on them. Denote
$a$, $b$ the corresponding correspondences from $\Spec(k)$ to
$X$, $Y$. Then $\alpha \times \beta$ corresponds to the
correspondence $a \otimes b$ from $\Spec(k)$ to $X \otimes Y = X \times Y$.
Hence the requirement follows from the fact that $G$ is
compatible with the tensor structures on both sides.

\medskip\noindent
Axiom (C)(d) follows because the cycle $[\Spec(k)]$
corresponds to the identity morphism on $h(\Spec(k))$.
This finishes the proof of the lemma.
\end{proof}

\end{document}
